\section{New classical families from balanced \texorpdfstring{\(\Ein\)}{Ein} coordinates}
\label{sec:new-classical-families}

The pure-exponential base-change theorem turns every full-rank linear image
of distinct \(\Ein\)-coordinates into an exact algebraic-independence
statement.  Two consequences are particularly concrete: a pairwise
algebraicity exclusion at every positive algebraic level, and a countable
family that is simultaneously algebraically independent and sharply small.

\subsection{Four classical values at every algebraic level}

\begin{theorem}[Four classical values at an algebraic level]
\label{thm:four-classical-algebraic-level}
For every positive algebraic number \(a\), among
\begin{equation}\label{eq:four-classical-level-values}
 E_1(a),\qquad \Ei(a),\qquad \Ci(a),\qquad \Chi(a)
\end{equation}
at most one is algebraic.

More precisely, each of the following six quantities is transcendental over
\(\Eexp\):
\begin{align}
 E_1(a)+\Ei(a),& & E_1(a)+\Ci(a),& & E_1(a)+\Chi(a),\notag\\
 \Ei(a)-\Ci(a),& & \Ei(a)-\Chi(a),& & \Chi(a)-\Ci(a).
 \label{eq:four-classical-witnesses}
\end{align}
\end{theorem}

\begin{proof}
Put
\[
 A=\Ein(a),\qquad B=\Ein(-a),\qquad
 C=\frac{\Ein(ia)+\Ein(-ia)}2.
\]
The connection formulas give, in the order displayed in
\eqref{eq:four-classical-witnesses},
\begin{equation}\label{eq:four-classical-witness-identities}
 A-B,\qquad A-C,\qquad \frac{A-B}{2},\qquad
 C-B,\qquad \frac{A-B}{2},\qquad
 C-\frac{A+B}{2}.
\end{equation}
The four values
\(\Ein(a),\Ein(-a),\Ein(ia),\Ein(-ia)\) are algebraically independent
over \(\Eexp\).  Every expression in
\eqref{eq:four-classical-witness-identities} is a nonzero linear form in
those independent coordinates and is therefore transcendental over
\(\Eexp\).  If any two members of
\eqref{eq:four-classical-level-values} were algebraic, the corresponding
witness in \eqref{eq:four-classical-witnesses} would be algebraic, a
contradiction.
\end{proof}

\subsection{A sharply decaying algebraically independent family}

\begin{theorem}[Quadratic balanced \(E_1\) family]
\label{thm:quadratic-balanced-E1}
For \(m\geq1\), put
\begin{equation}\label{eq:quadratic-balanced-E1-definition}
 Q_m=E_1(m^2)-2E_1\bigl(m(m+1)\bigr)+E_1\bigl((m+1)^2\bigr).
\end{equation}
Then the countable family \(\{Q_m:m\geq1\}\) is algebraically independent
over \(\Eexp\), in the finite-subset sense.  Every \(Q_m\) is positive, and
\begin{equation}\label{eq:quadratic-balanced-E1-asymptotic}
 Q_m=\frac{e^{-m^2}}{m^2}
 \left(1+O(m^{-2})+O(e^{-m})\right).
\end{equation}
Consequently
\begin{equation}\label{eq:quadratic-balanced-E1-saturation}
 \boxed{
 \lim_{m\to\infty}
 \frac{-\log Q_m}{(m+1)^2}=1.}
\end{equation}
\end{theorem}

\begin{proof}
The two cancellations
\[
 1-2+1=0,
 \qquad
 \log(m^2)-2\log\bigl(m(m+1)\bigr)+\log\bigl((m+1)^2\bigr)=0
\]
give the exact identity
\begin{equation}\label{eq:quadratic-balanced-Ein-form}
 Q_m=\Ein(m^2)-2\Ein\bigl(m(m+1)\bigr)
       +\Ein\bigl((m+1)^2\bigr).
\end{equation}
For \(m\geq1\), the integer \(m(m+1)\) is not a square: since
\(\gcd(m,m+1)=1\), a square product would force both factors to be
squares, and two consecutive positive squares do not exist.  Moreover
\(m(m+1)\) is strictly increasing with \(m\).  Hence the middle coordinate
\(\Ein(m(m+1))\) occurs in no row of
\eqref{eq:quadratic-balanced-Ein-form} other than the \(m\)-th row.  In any
finite linear combination of these rows, inspection of the middle
coordinates forces every row coefficient to vanish.  The rows therefore
have full rank.  The pure-exponential base-change theorem and
Lemma~\ref{lem:linear-image-ein} prove the asserted algebraic independence.

For \(x>0\), integration by parts gives the standard bounds
\begin{equation}\label{eq:E1-two-sided-bound}
 \frac{e^{-x}}{x+1}<E_1(x)<\frac{e^{-x}}x.
\end{equation}
Indeed, the upper bound is immediate, while
\[
 E_1(x)=\frac{e^{-x}}x-\int_x^\infty\frac{e^{-t}}{t^2}\,dt
 >\frac{e^{-x}}x-\frac{E_1(x)}x
\]
gives the lower bound.  Therefore
\[
 \frac{2E_1(m^2+m)}{E_1(m^2)}
 <2e^{-m}\frac{m^2+1}{m^2+m}
 \leq\frac2e<1.
\]
Thus \(E_1(m^2)>2E_1(m^2+m)\), and the last positive term in
\eqref{eq:quadratic-balanced-E1-definition} proves \(Q_m>0\).

Finally, the standard expansion
\[
 E_1(x)=\frac{e^{-x}}x\left(1+O(x^{-1})\right)
\]
shows that
\[
 \frac{E_1(m^2+m)}{E_1(m^2)}=O(e^{-m}),
 \qquad
 \frac{E_1((m+1)^2)}{E_1(m^2)}=O(e^{-2m}).
\]
This proves \eqref{eq:quadratic-balanced-E1-asymptotic}; taking logarithms
gives \eqref{eq:quadratic-balanced-E1-saturation}.
\end{proof}

\begin{remark}[Uniform moving-point boundary]
The scale \((m+1)^2\) in
\eqref{eq:quadratic-balanced-E1-saturation} is the largest evaluation
argument, not a Weil height or a relation-polynomial height.  The theorem
therefore gives a sharp counterexample to naive lower bounds that are
uniform only in the moving archimedean evaluation scale.  It does not
contradict fixed-point algebraic-independence measures for \(E\)-functions.
\end{remark}

\subsection{A conditional \texorpdfstring{\(E\)--\(D\)}{E--D} collision}

\begin{corollary}[An explicit conditional \(E\)--\(D\) collision]
\label{cor:gamma-algebraic-ED-collision}
Let \(\mathbf E\) and \(\mathbf D\) be the rings of finite \(E\)-values
and arithmetic-Gevrey values of order one, respectively, in the sense of
Fischler--Rivoal \cite{FR2024ED}.  Then
\begin{equation}\label{eq:gamma-algebraic-ED-collision}
 \boxed{
 \gamma\in\Qbar
 \quad\Longrightarrow\quad
 \delta\in(\mathbf E\cap\mathbf D)\setminus\Qbar.}
\end{equation}
Consequently the conjectural equality
\(\mathbf E\cap\mathbf D=\Qbar\) implies that \(\gamma\) is
transcendental.
\end{corollary}

\begin{proof}
Assume \(\gamma\in\Qbar\).  Then
\[
 f(z)=e^z\bigl(\Ein(z)-\gamma\bigr)
\]
is an \(E\)-function with algebraic Taylor coefficients, and
\(f(1)=\delta\).  Hence \(\delta\in\mathbf E\).  The standard Borel-sum
realization of Euler's factorial series gives \(\delta\in\mathbf D\).
If \(\delta\) were algebraic, then
\[
 \Ein(1)=\gamma+e^{-1}\delta
\]
would belong to \(\Eexp\), contradicting the pure-exponential base-change
theorem.  Thus \(\delta\notin\Qbar\), proving the assertion.
\end{proof}
